\section{8차시 --- 알비 카트 주행}
\label{sec:ch08}

본 차시는 7차시에서 익힌 DC모터에 바퀴를 결합하여 "움직이는 로봇"을
완성하는 단계이다. 전진/후진/속도 조절을 결합해 단순한 주행에서
왕복 주행, 부드러운 가감속, 그리고 학습자 자유 코스로 발전시킨다.

\subsection{미션 1 --- 알비 카트 앞으로!}
\label{subsec:ch08-m1}
\missionmeta{DC모터, 바퀴 2개}{MOTOR}

\begin{storybox}
\sayares{알비! 바퀴를 달면 앞으로 달릴 수 있어?}
\sayalbi{삐릿- 물론이에요! DC모터가 정방향으로 돌면 바퀴가 구르면서 앞으로 나아가요!}
\end{storybox}

\subsubsection*{학습 목표}
\begin{itemize}[leftmargin=1.4em]
  \item DC모터 정방향으로 회전 → 전진
  \item 3초 전진 후 정지
  \item 알비 카트가 앞으로 달리는 것 확인
\end{itemize}

\begin{labbox}{샘플 코드 --- 알비 카트 전진!}
\begin{minted}{python}
# 알비 카트 전진!
motor_forward(3)   # 3초 전진
motor_stop()
\end{minted}
\end{labbox}

\subsection{미션 2 --- 알비 카트 뒤로! 왕복 주행}
\label{subsec:ch08-m2}
\missionmeta{DC모터, 바퀴 2개}{MOTOR}

\begin{storybox}
\sayalbi{삐릿! 아레스, 이번엔 뒤로도 달려봐요! DC모터 방향을 반대로 바꾸면 후진!}
\sayares{앞으로 갔다가 뒤로 오면 왕복 주행이 되겠네! 해보자!}
\end{storybox}

\subsubsection*{학습 목표}
\begin{itemize}[leftmargin=1.4em]
  \item DC모터 역방향 회전 → 후진
  \item 전진(3초) → 정지(1초) → 후진(3초) → 정지
  \item 왕복 주행 완성
\end{itemize}

\begin{labbox}{샘플 코드 --- 왕복 주행}
\begin{minted}{python}
# 왕복 주행
motor_forward(3)    # 전진
motor_stop()
time.sleep(1)
motor_backward(3)   # 후진
motor_stop()
\end{minted}
\end{labbox}

\subsection{미션 3 --- 속도 조절 주행!}
\label{subsec:ch08-m3}
\missionmeta{DC모터, 바퀴 2개}{MOTOR}

\begin{storybox}
\sayares{알비, 항상 같은 속도로만 달리면 재미없어. 천천히 출발해서 점점 빨라지다가 멈추는 건?}
\sayalbi{삐릿- 좋아요! PWM으로 속도를 조절하면 부드럽게 달릴 수 있어요!}
\end{storybox}

\subsubsection*{학습 목표}
\begin{itemize}[leftmargin=1.4em]
  \item 느리게 출발: 30\% 속도로 2초
  \item 중간 속도: 60\% 속도로 2초
  \item 전속력: 100\% 속도로 2초
  \item 서서히 감속 70→40→10→정지 부드러운 정차
\end{itemize}

\begin{labbox}{샘플 코드 --- 속도 조절 주행}
\begin{minted}{python}
# 속도 조절 주행
motor_forward_pwm(30)
time.sleep(2)
motor_forward_pwm(60)
time.sleep(2)
motor_forward_pwm(100)
time.sleep(2)
motor_forward_pwm(70); time.sleep(0.5)
motor_forward_pwm(40); time.sleep(0.5)
motor_forward_pwm(10); time.sleep(0.5)
motor_stop()
\end{minted}
\end{labbox}

\subsection{미션 4 --- 알비 카트 나만의 코스!}
\label{subsec:ch08-m4}
\missionmeta{DC모터, 바퀴 2개}{MOTOR}

\begin{storybox}
\sayares{알비! 전진, 후진, 속도 조절 다 알았으니 나만의 주행 코스를 만들어볼게!}
\sayalbi{삐릿♪ 전진했다가 멈추고 후진하는 나만의 드라이브 코스를 코딩해봐요!}
\end{storybox}

\subsubsection*{학습 목표}
\begin{itemize}[leftmargin=1.4em]
  \item 전진(2초) → 정지(0.5초) → 후진(2초) → 정지 반복
  \item 속도 변화로 역동적 주행 연출
  \item 나만의 알비 카트 코스 완성 + 시연
\end{itemize}

\begin{labbox}{샘플 코드 --- 나만의 주행 코스}
\begin{minted}{python}
# 나만의 주행 코스
for i in range(3):
    motor_forward(2)
    motor_stop()
    time.sleep(0.5)
    motor_backward(2)
    motor_stop()
    time.sleep(0.5)
\end{minted}
\end{labbox}

\begin{summary}[8차시 요약]
7차시에서 익힌 모터 제어가 "바퀴"라는 기계 부품과 만나면서 비로소
"로봇 주행"이라는 통합 시스템이 완성된다. 학습자는 PWM 듀티 사이클의
연속적 변화가 카트의 가속/감속 곡선으로 그대로 매핑됨을 체험적으로
이해한다.
\end{summary}
